% Section 5: the C implementation and what it checks. The numbers it produced
% are in Section 6, where they are compared rather than merely listed.
\section{Implementation in C}
\label{sec:implementation}

The subject leaves code optional and provides a skeleton in C. It is written
here because several claims in this report are the kind that should be measured
rather than asserted: that generation always returns a spanning tree, that both
outputs return the same route, and that reading the tree is cheaper than
searching the graph.

\subsection{What was built}

C11, one Makefile, no dependency outside the standard library, compiled with
\texttt{-Wall -Wextra -Wpedantic -Werror}. It lives beside this report in
\texttt{implementation/}.

\begin{center}
\begin{tabular}{@{}ll@{}}
  \toprule
  File & Role \\
  \midrule
  \texttt{cell.\{h,c\}}        & a cell, and the two axis operations of
                                 Section~\ref{sec:regions} \\
  \texttt{maze.\{h,c\}}        & the first output: nodes, coordinates, neighbours \\
  \texttt{region\_tree.\{h,c\}} & the second output: divisions, doors, addresses \\
  \texttt{generator.\{h,c\}}   & Algorithm~\ref{alg:generate-maze}, filling both \\
  \texttt{solver.\{h,c\}}      & Algorithm~\ref{alg:search} on the graph,
                                 Algorithm~\ref{alg:region-route} on the tree \\
  \texttt{invariants.\{h,c\}}  & the properties a finished maze must have \\
  \texttt{svg.\{h,c\}}         & drawing a maze, and a route through it \\
  \texttt{tests.c}             & the checks below \\
  \texttt{benchmark.c}         & the measurements of Section~\ref{sec:analysis} \\
  \bottomrule
\end{tabular}
\end{center}

Two decisions are worth naming. The invariants are their own module, because a
generator that checked its own work would only ever confirm its own
assumptions. And the random numbers come from a small \texttt{xorshift}
generator rather than \texttt{rand()}, so that a seed determines a maze on any
machine and the measurements are reproducible.

\subsection{What is checked}

$1121$ assertions, all passing.

\begin{description}[leftmargin=*,style=nextline,itemsep=4pt]
  \item[Generation returns a spanning tree.] $WH-1$ doors, connected, no cell
    with more than four, over ten shapes $\times$ twenty seeds including
    $1\times1$, $1\times40$ and $40\times1$.
  \item[A corridor is laid down whole.] A $1\times40$ area gives a tree of one
    leaf and $39$ doors --- Section~\ref{sec:corridor-base} observed rather than
    argued.
  \item[The same seed gives the same maze.] Without which no measurement below
    is reproducible.
  \item[Both outputs agree, cell for cell.] The graph search and the tree search
    return identical routes on every shape and seed tried. This is the claim the
    whole of Section~\ref{sec:common-ancestor} rests on.
  \item[Every guide returns the same route.] Plain Dijkstra and all four
    heuristic settings agree, which is
    Proposition~\ref{prop:spanning-tree} made visible.
  \item[The route is walkable.] Consecutive cells on it are joined by a door.
  \item[Bad input is refused.] A zero dimension, a door between cells that do
    not touch, a door opened twice, a cell outside the grid.
\end{description}

\begin{keypoint}
The fourth item is the one that matters. Two algorithms that share no code ---
one walking the graph with a priority queue, one descending the tree with
rectangle tests --- return \keyterm{the same cells in the same order}, every
time.
\end{keypoint}
